\expandafter\def\csname HistoricalNoteText2\endcsname{It should be noted that, in his own account of this proof, the distinguished \textsc{d'Alembert} used geometrical considerations, regarding $X$ as the abscissa and $x$ as the ordinate of a curve (following the practice of all mathematicians in the first part of this century, among whom the notion of functions was less familiar). Since, however, all his reasoning, considered only in its substance, rests on no geometrical principles but on purely analytical ones, and the expressions imaginary curves and imaginary ordinates seem rather harsh and more liable to offend the present-day reader, I have preferred a purely analytical presentation here. I added this note so that no one comparing \textsc{d'Alembert}'s own proof with this brief account should suspect that anything essential has been changed.}

\expandafter\def\csname HistoricalNoteText3\endcsname{I note in passing on this occasion that many of these series seem at first sight extremely convergent: for example, most of those that the distinguished \textsc{Euler} uses in the latter part of the \textit{Inst. Calc. Diff.}, Chapter VI, pp. 441--474, to assign as closely as possible the sums of other series (the remaining series on pp. 475--478 can indeed converge). As far as I know, no one has observed this before. It would therefore be very desirable to show clearly and rigorously why series of this kind, which at first converge very rapidly, then gradually more and more slowly, and finally diverge more and more, nevertheless give a nearly true sum provided not too many terms are taken, and how far such a sum can safely be regarded as exact?}

